\section{Counting connected clusters in a graph of bounded degree}
\label{sec:hastings_connected_set_count}

% Hastings, arXiv:cond-mat/0508554v3, proj3new.tex, lines 1227--1254 (ending at
% eq. perc2); the classical count of lattice animals is Klarner, Canad. J. Math. 19
% (1967) 851.

The high-temperature cluster expansion of \cite{Hastings2006Solving} sums over
clusters of sites, and its convergence rests on the classical fact that a
lattice of $N$ sites carries at most $Nx^{n}$ clusters of $n$ sites for some
constant $x$ \cite[lines 1227--1254]{Hastings2006Solving}. We prove this for an
arbitrary finite simple graph whose vertices all have degree at most $\Delta$,
with the explicit constant $x=\max(1,\Delta^{2})$. A cluster is a set of vertices that
induces a connected subgraph.

\begin{theorem}[Clusters through a vertex]
    \label{thm:hastings_connected_set_count_mem}
    \lean{TNLean.Lattice.card_connected_induce_mem_le}
    \leanok
    Let $G$ be a finite simple graph in which every vertex has degree at most
    $\Delta$, let $v$ be a vertex, and let $n$ be a natural number. At most
    $\Delta^{2(n-1)}$ sets of $n$ vertices contain $v$ and induce a connected
    subgraph of $G$, where $n-1$ is read as $0$ when $n=0$; no such set
    exists then.
\end{theorem}

\begin{proof}
    \leanok
    Let $S$ be such a set. Starting from $\{v\}$ and the walk of length zero,
    add the vertices of $S$ one at a time: as long as the set $T$ reached so
    far is a proper subset of $S$, a path inside $S$ from $v$ to a vertex of
    $S\setminus T$ leaves $T$ along an edge $wu$ with $w\in T$ and
    $u\in S\setminus T$. Inserting the detour $w\to u\to w$ into the current
    closed walk at a visit of $w$ adds two to its length and exactly the vertex
    $u$ to its vertex set. After $n-1$ steps this gives a closed walk from $v$
    of length $2(n-1)$ whose vertex set is $S$. The walk determines $S$, so
    distinct sets give distinct walks, and at most $\Delta^{L}$ walks of length
    $L$ start at $v$, since each step has at most $\Delta$ choices.
\end{proof}

\begin{theorem}[Clusters of a given size]
    \label{thm:hastings_connected_set_count}
    \lean{TNLean.Lattice.card_connected_induce_le,
        TNLean.Lattice.card_connected_induce_le_exponential}
    \leanok
    \uses{thm:hastings_connected_set_count_mem}
    Let $G$ be a simple graph on $N$ vertices in which every vertex has degree
    at most $\Delta$, and let $n$ be a natural number. At most
    $N\Delta^{2(n-1)}$ sets of $n$ vertices induce a connected subgraph of $G$,
    with $n-1$ read as $0$ when $n=0$.
    In particular the number is at most $Nx^n$ for
    $x=\max(1,\Delta^2)$. The latter bound also holds for $n=0$, since the
    empty induced graph is not connected.
\end{theorem}

\begin{proof}
    \leanok
    \uses{thm:hastings_connected_set_count_mem}
    A set inducing a connected subgraph is nonempty, so it contains some
    vertex; sum the bound of Theorem~\ref{thm:hastings_connected_set_count_mem}
    over the $N$ vertices.
    For $n\ge1$, $\Delta^{2(n-1)}=(\Delta^2)^{n-1}\le x^{n-1}\le x^n$,
    since $x\ge1$. This choice also covers $\Delta=0$, when the connected
    singleton sets would prevent taking $x=0$.
\end{proof}
